\section*{الفصل \ref{ch:b3:nervous-systems} --- تنظيم الأجهزة العصبية}

\begin{solution}{exo:b3:nervous-systems:1}
الخلايا النجمية: توازن البوتاسيوم والغلوتامات، وتغذي
العصبونات، وتستحث \omterm{prop:b3:nervous-systems:barrier-budget}{الحاجز الدموي الدماغي}، وتنظم الجريان
الدموي الموضعي. وقليلات التغصن: تميلن المحاوير في الجهاز
العصبي المركزي. وخلايا شوان: تميلن المحاوير المحيطية (وتوجّه
تجددها). والدبقيات الصغيرة: بلعميات الدماغ، تقلّم المشابك
وتصفّي الحطام. (وتبطّن الخلايا البِطانية البُطينات وتصنع
\omterm{prop:b3:nervous-systems:barrier-budget}{السائل الدماغي الشوكي}.)
\end{solution}

\begin{solution}{exo:b3:nervous-systems:2}
سوّدت الصبغة عصبونات قليلة تسويدًا كاملًا وتركت الباقي غير
مرئي، فأمكن أن تُرى خلية واحدة كاملة. واستنتج كاخال أن
العصبونات خلايا منفصلة يلامس بعضها بعضًا من غير أن تندمج،
وأن للجريان اتجاهًا ثابتًا من التغصنات إلى المحوار؛ واستنتج
غولجي أنها تكوّن شبكة متصلة. وكان كاخال محقًا؛ وأظهر المجهر
الإلكتروني الفجوة المشبكية عام 1954.
\end{solution}

\begin{solution}{exo:b3:nervous-systems:3}
طرقة تمطّط رباعية الرؤوس؛ فتشتعل مغازلها العضلية؛ وتدخل
واردات Ia إلى النخاع وتتشابك مباشرةً على العصبونات الحركية
لرباعية الرؤوس، فتشتعل وتقبض العضلة، بينما يثبّط \omterm{def:b3:nervous-systems:cells}{عصبون بيني}
العصبونات الحركية لأوتار المأبض. فمشبك واحد، ومحاوير قصيرة
عند \qty{80}{m/s}: أي نحو \qty{30}{ms}. وأما الركلة الإرادية
فتمر بالدماغ --- إدراك، وقرار، وتخطيط حركي، وسبل نازلة،
ومشابك عدة --- فتستغرق \qtyrange{200}{300}{ms}.
\end{solution}

\begin{solution}{exo:b3:nervous-systems:4}
الودّي (نورأدرينالين عند الهدف): قلب أسرع وأقوى، وحدقة
متسعة، وحركية الأمعاء وإفرازها منخفضان، وطرق هوائية متوسعة.
ونظير الودي (أسيتيل كولين): قلب أبطأ، وحدقة متضيقة، وحركية
الأمعاء وإفرازها مرتفعان، وطرق هوائية متضيقة.
\end{solution}

\begin{solution}{exo:b3:nervous-systems:5}
$\lambda = \sqrt{R_{m}d/4R_{i}} = \sqrt{10^{4}\times 10^{-4}/400}$ سنتيمتر
$= \qty{0.05}{cm} = \qty{0.5}{mm}$. وبعد \qty{2}{mm}: $e^{-4} \approx
2\,\%$. ومضاعفة $\lambda$ تحتاج إلى أربعة أمثال القطر، أي
\qty{4}{\micro m}.
\end{solution}

\begin{solution}{exo:b3:nervous-systems:6}
الحسي: $12\times 6 = \qty{72}{m/s}$، أي $1.2/72 = \qty{17}{ms}$.
والحركي: \qty{90}{m/s}، أي $1.1/90 = \qty{12}{ms}$. ومشبكان،
\qty{1}{ms}: فالمجموع نحو \qty{30}{ms}. وغير المميلن:
$1.1\sqrt{12} = \qty{3.8}{m/s}$ و$1.1\sqrt{15} = \qty{4.3}{m/s}$: أي
$315 + 258 + 1 \approx \qty{570}{ms}$ --- وهو أبطأ من أن ينقذ
القدم.
\end{solution}

\begin{solution}{exo:b3:nervous-systems:7}
$2.4\times 10^{20}/8.6\times 10^{10} \approx 2.8\times 10^{9}$ من ATP
لكل عصبون في الثانية؛ وعند $2\times 10^{9}$ لكل شوكة، يكون
المعدل نحو \qty{1.4}{Hz} وسطيًا. فالترميز لا بد أن يكون
متناثرًا: شوكات قليلة، والمعلومة محمولة بأي العصبونات تشتعل
ومتى، لا بمعدلات عالية في كل مكان.
\end{solution}

\begin{solution}{exo:b3:nervous-systems:8}
عصبونان يثبّط كل منهما الآخر بلا تعب يكوّنان قاطعة: فأولهما
اشتعالًا يكبت الآخر ما دام الدفع قائمًا، ولا يقع تسليم أبدًا.
والتذبذب يحتاج إلى عملية بطيئة تُضعف قبضة الفائز ---
تكيّفًا في اشتعاله، أو اكتئابًا في مشبكه المثبِّط، أو ارتدادًا
للخاسر بعد تحرره --- ويحدد ثابتُ زمنها الدورَ.
\end{solution}

\begin{solution}{exo:b3:nervous-systems:9}
يُحمل L-dopa عبر \omterm{def:b3:innate-immunity:innate}{الحاجز} بناقل الأحماض الأمينية المحايدة
الكبيرة؛ وأما الدوبامين، وهو قطبي مشحون، فلا. فيتناول
المرضى L-dopa (مع مثبط للإنزيم المحيطي حتى لا يتحول قبل
العبور)، ويصنع إنزيم الدماغ نفسه الدوبامين منه حيث يلزم.
وأما الأضداد، عند \qty{150}{kDa}، فلا تعبر الوصلات المحكمة
ولا أي ناقل؛ ويحتاج إيصالها إلى مكوك يختطف مستقبلًا للنقل
عبر الخلية، أو إلى حقن في \omterm{prop:b3:nervous-systems:barrier-budget}{السائل الدماغي الشوكي}.
\end{solution}

\begin{solution}{exo:b3:nervous-systems:10}
تيار الغشاء في وحدة الطول هو المقاوم $V/r_{m}$ مضافًا إليه
السعوي $c_{m}\,\partial V/\partial t$؛ ويعطي حفظ الشحنة
$\partial I/\partial x = -(V/r_{m} + c_{m}\,\partial V/\partial t)$،
ومع $\partial V/\partial x = -r_{i}I$ نجد $\frac{1}{r_{i}}\,\partial^{2}V/
\partial x^{2} = V/r_{m} + c_{m}\,\partial V/\partial t$؛ وبالضرب في
$r_{m}$: $\lambda^{2}\,\partial^{2}V/\partial x^{2} = V + \tau\,\partial
V/\partial t$. وعند $\partial V/\partial t = 0$ تعود معادلة
المبرهنة. و$\lambda$ طول و$\tau$ زمن، فتكون $\lambda/\tau$
سرعة --- وهي المقياس الطبيعي لسرعة انتشار اضطراب.
\end{solution}

\begin{solution}{exo:b3:nervous-systems:11}
$v \propto \sqrt{d}$: فأربعة أمثال السرعة تحتاج إلى ستة عشر
مثلًا من القطر، أي \qty{8}{mm}. والمساحة $\pi(4\,\text{mm})^{2} \approx \qty{50}{mm^2}$
في مقابل $\pi\,(10\,\unit{\micro m})^{2} = \qty{314}{\micro m^2}$: أي نسبة
$1.6\times 10^{5}$. فعصب من ألياف سريعة غير مميلنة مستحيل ---
إذ يكون ليف واحد منها بسماكة العصب كله --- فالحيوان ذو
القنوات السريعة الكثيرة يحتاج إلى \omterm{def:b3:nervous-systems:cells}{المييلين}، ولذلك تطور عند
الفقاريات (وعند بضعة لافقاريات استقلالًا).
\end{solution}

\begin{solution}{exo:b3:nervous-systems:12}
$a = 0.4/30^{0.75} = 0.4/12.8 = 0.031$؛ وعند كتلة \qty{70000}{g}: $0.031\times
\num{70000}^{0.75} = 0.031\times 4300 \approx \qty{134}{g}$. ودماغ الإنسان
\qty{1350}{g} عشرة أمثال التنبؤ. وتتبع كتلة الدماغ كتلةَ
الجسم لأن الجسم الأكبر يحتاج إلى عصبونات حسية وحركية أكثر،
ولأن العصبونات نفسها تكبر مع حجم الدماغ؛ والذي يتابع
الإدراك هو عدد العصبونات في القشرة، وهو يتوقف على كثافة
رصّها --- والرئيسيات ترصّ في الغرام أكثر من سائر الثدييات،
والإنسان أكثرها جميعًا.
\end{solution}

\begin{solution}{pb:b3:nervous-systems:1}
\textbf{1.} $\lambda = \sqrt{2\times 10^{4}\,d/400}$: فعند $d =
\qty{1e-4}{cm}$، $\sqrt{5\times 10^{-3}} = \qty{0.071}{cm} =
\qty{0.71}{mm}$؛ وعند \qty{4}{\micro m}، \qty{1.41}{mm}. و$\tau = 2\times
10^{4}\times 10^{-6} = \qty{20}{ms}$.
\textbf{2.} $5\,e^{-0.3/0.71} = \qty{3.3}{mV}$؛ و$5\,e^{-0.3/1.41} =
\qty{4.0}{mV}$.
\textbf{3.} الدخلات البعيدة تصل موهنة، فيحدد موضع وقوع المشبك
وزنَه؛ والكمون يهبط في نحو $\tau$، فلا يتجامع دخلان إلا إذا
وقعا في نحو \qty{20}{ms} أحدهما من الآخر --- فالعصبون كاشف
تزامن بنافذة \qty{20}{ms}.
\textbf{4.} المميلن \qty{60}{m/s}؛ وغير المميلن $1.1\sqrt{10} =
\qty{3.5}{m/s}$؛ ولبلوغ \qty{60}{m/s} من دون \omterm{def:b3:nervous-systems:cells}{مييلين} يلزم $d = (60/1.1)^{2}
\approx \qty{3000}{\micro m}$: أي ثلاثة ميليمترات.
\textbf{5.} $\pi(5)^{2} = \qty{79}{\micro m^2}$ في مقابل $\pi(1500)^{2} =
7\times 10^{6}\,\unit{\micro m^2}$: أي نحو \num{90000} محوار مميلن
في حيز محوار واحد.
\textbf{6.} \qty{2}{cm} عند \qty{60}{m/s} هي \qty{0.3}{ms}؛ وعند
\qty{3.5}{m/s}، \qty{5.7}{ms}: أي نحو \qty{5}{ms} من التأخير
الإضافي. والأسوأ أن للغشاء العاري سعة كبيرة وقنوات قليلة،
فقد يخفق التيار الآتي من آخر عقدة سليمة في بلوغه العتبة:
أي حصار النقل.
\textbf{7.} الحسي $1.2/60 = \qty{20}{ms}$؛ ومشبكان \qty{1}{ms}؛
والحركي $1.1/90 = \qty{12}{ms}$: أي نحو \qty{33}{ms} (مع ميلي
ثانية عند الوصل العصبي العضلي).
\textbf{8.} إلى الدماغ: $0.6/60 = \qty{10}{ms}$ مع مشبكين، أي
\qty{11}{ms} بعد دخول الإشارة إلى النخاع عند \qty{20}{ms}:
فتبلغ الإشارة القشرة في نحو وقت تحرك القدم؛ وتحتاج الخبرة
الواعية بالألم إلى بضع مئات من الميلي ثانية أخرى من
المعالجة، فتُسحب القدم قبل أن يُحسّ الألم.
\textbf{9.} $1.1\sqrt{1} = \qty{1.1}{m/s}$: أي $1.8/1.1 \approx
\qty{1.6}{s}$. فالألم الأول الحاد يصل بألياف مميلنة رفيعة في
عشرات الميلي ثانية؛ والألم الثاني الكليل الحارق بألياف C
غير المميلنة بعد ثانية أو أكثر.
\textbf{10.} الحسي $3.2/60 = \qty{53}{ms}$، والحركي $3.1/90 =
\qty{34}{ms}$، والمشابك \qty{1}{ms}: أي نحو \qty{90}{ms}.
فالحيوانات الطويلة منعكساتها بطيئة؛ وهي تعوّض ذلك بمحاوير
أسمك وبتفويض أكثر إلى الدارات الموضعية.
\textbf{11.} $0.8/80 = \qty{10}{ms}$ في كل اتجاه، أي \qty{20}{ms}؛
والمشبك \qty{0.5}{ms}، والوصل \qty{1}{ms}، والقوة \qty{5}{ms}:
أي \qty{26.5}{ms}، ويكمل تحويلُ المغزل نفسه والنقلُ داخل
العضلة بقية \qty{30}{ms}.
\textbf{12.} المحاوير غير المميلنة تنقل بمتر أو مترين في
الثانية، فالعُرى التي تستغرق \qty{30}{ms} في البالغ تستغرق
مئات، بتوقيت رديء؛ والتميلين يرفع السرعة عشرة إلى خمسين
ضعفًا ويجعل التوقيت دقيقًا، وهو ما يتطلبه المشي والقبض
والكلام.
\textbf{13.} $20/50\,000 = 4\times 10^{-4}$ مول في الثانية $= 2.4\times 10^{20}$
من ATP في الثانية؛ أي $2.8\times 10^{9}$ لكل عصبون في الثانية.
\textbf{14.} $4\times 10^{8} + 7000\times 2\times 10^{4} = 4\times 10^{8}
+ 1.4\times 10^{8} = 5.4\times 10^{8}$ من ATP.
\textbf{15.} إذا ذهبت كلها إلى الشوكات: $2.8\times 10^{9}/5.4\times 10^{8} \approx
\qty{5}{Hz}$؛ وإذا ذهب نصفها: نحو \qty{2.6}{Hz}.
\textbf{16.} عند \qty{1}{Hz}، $5.4\times 10^{8}/2.8\times 10^{9} \approx
\qty{20}{\%}$ من الميزانية: فالترميز متناثر --- معظم
العصبونات صامتة معظم الوقت، والمعلومة في أيها القليل يشتعل
ومتى.
\textbf{17.} $8.6\times 10^{10}\times 50\times 5.4\times 10^{8} = 2.3
\times 10^{21}$ من ATP في الثانية $= 3.9\times 10^{-3}$ مول في الثانية $\approx \qty{190}{W}$
--- أي ضعف الجسم كله في الراحة.
\textbf{18.} $5\times 16 = \qty{80}{kJ/h} = \qty{22}{W}$: أي يكفي
بالكاد، بلا احتياطي. وحين ينقطع الإمداد تتوقف المضخات في
ثوانٍ، وتستنفد التدرجات، ويُفقد الوعي في نحو عشر ثوانٍ.
\textbf{19.} $a = 0.4/12.8 = 0.031$؛ وعند \qty{70}{kg}: $0.031\times 4300 =
\qty{134}{g}$؛ وعند \qty{5000}{kg}: $0.031\times 1.06\times 10^{5} \approx
\qty{3300}{g}$.
\textbf{20.} الإنسان $1350/134 \approx 10$؛ والفيل $4800/3300 \approx
1.5$.
\textbf{21.} في \omterm{def:b3:nervous-systems:plan}{المخيخ} --- \qty{98}{\%} منها --- تنسّق أربعين
ألف عضلة في جذعه. فعدد العصبونات القشرية، لا المجموع، هو
الذي يتابع الإدراك؛ \omterm{def:b3:nervous-systems:plan}{والمخيخ} يتابع المهمة الحركية في إدارة
جسم ضخم.
\textbf{22.} تطول المحاوير والتغصنات مع الدماغ، فيشغل كل
عصبون حجمًا أكبر وينمو عدد العصبونات أبطأ من الكتلة. وفي
الرئيسيات يبقى حجم العصبون شبه ثابت مع نمو الدماغ، فينمو
عدد العصبونات على التناسب مع الكتلة تقريبًا، ويحمل دماغ
رئيسي بوزن معين أضعاف ما يحمله دماغ قارض بذلك الوزن.
\textbf{23.} تقدر الذراع على تنسيق ممصّاتها ومفاصلها، وعلى
المدّ، والقبض، وتمرير الطعام على نفسها، وتجنّب قبض جلدها هي،
بدارات موضعية؛ ولا تقدر على أن ترى، ولا أن تختار هدفًا، ولا
أن تتعلم تمييزًا. والاختبار: اقطع حبل الذراع العصبي عن
الدماغ والمس الذراع بطعام --- فتقبضه وتمرره نحو حيث يكون
الفم؛ وتفعل ذراع مبتورة الشيء نفسه ساعةً.
\textbf{24.} نحو $23$ مشبكًا لكل عصبون في مقابل $7000$: أي أقل
بثلاثمئة ضعف. والخريطة الكاملة تقول أي عصبون يقدر على
التأثير في أي عصبون، لا قوة كل وصلة ولا إشارتها، ولا أي
المعدّلات العصبية تغيّرها، ولا الحركيات --- فحتى في الدودة
لا يتنبأ مخطط التوصيل بالسلوك إلا جزئيًا.
\textbf{25.} $\lambda = \qty{0.71}{mm}$ للتغصن قطره
\qty{1}{\micro m}؛ وكمون السحب نحو \qty{33}{ms}؛ ومعدل
الاشتعال الوسطي نحو \qty{2.6}{Hz} مع نصف الميزانية على
الشوكات.
\end{solution}
